% !TEX root = ../thesis.tex
% ==============================================================================
%  CHAPTER 6: CONCLUSIONS & FUTURE WORK (فصل ششم: نتیجه‌گیری و کارهای آینده) 
%             (chapters/chapter6_conclusions.tex)
% ==============================================================================
%  This chapter summarizes the primary findings and contributions of the
%  thesis, addresses theoretical and computational limitations, and outlines
%  promising directions for future research.
% ==============================================================================


\chapter{نتیجه‌گیری و کارهای آتی}\label{chap:chapter6_conclusions}

فصل پایانی پایان‌نامه وظیفه‌ی جمع‌بندی منسجم یافته‌ها، بازخوانی دستاوردهای اصلی، تبیین صریح محدودیت‌ها و ترسیم افق‌های پژوهشی آینده را بر عهده دارد. در این فصل، نویسنده به مخاطب نشان می‌دهد که چگونه پرسش‌های اولیه‌ی پژوهش که در فصل اول مطرح شده بودند، در طول نوشتار با تکیه بر مبانی نظری، نقد پیشینه، توسعه‌ی روش‌شناسی و ارزیابی‌های تجربی پاسخ داده شده‌اند.


\section{جمع‌بندی دستاوردها و نتایج محوری}\label{sec:6-1}

در بخش جمع‌بندی، مهم‌ترین دستاوردهای علمی و عملیاتی پژوهش به‌صورت خلاصه و با استناد به شواهد ارائه‌شده در فصول قبل بازگو می‌شوند. این بخش نباید تکرار واژه‌به‌واژه‌ی چکیده باشد، بلکه باید نمایی کلان از جایگاه نتایج و اثرگذاری آن‌ها بر حوزه‌ی موضوعی پژوهش ارائه دهد. نویسنده در این بخش به‌طور مشخص بیان می‌کند که چارچوب پیشنهادی تا چه میزان موفق به بهبود معیارهای هدف نسبت به روش‌های خط‌پایه شده و کدام فرضیات اولیه‌ی پژوهش از طریق آزمون‌های تجربی مورد تأیید قرار گرفته‌اند.


\section{محدودیت‌های پژوهش}\label{sec:6-2}

بیان صادقانه و واقع‌بینانه‌ی محدودیت‌ها از شاخص‌های برجسته‌ی بلوغ علمی پژوهشگر است و به اعتبار سند می‌افزاید. هر پژوهشی با محدودیت‌هایی روبه‌رو است که شناخت دقیق آن‌ها مرزهای کاربرد روش را روشن می‌سازد. این محدودیت‌ها می‌توانند ناشی از موارد زیر باشند:
\begin{itemize}
    \item مفروضات نظری و شرایط مرزی حاکم بر مدل‌سازی ریاضی؛
    \vspace{-0.25cm}
    \item محدودیت‌های دسترسی به داده‌ها، ابزارهای اندازه‌گیری یا منابع پردازشی؛
    \vspace{-0.25cm}
    \item دامنه‌ی تعمیم‌پذیری نتایج به سناریوها و کاربردهای متفاوت بیرونی.
\end{itemize}


\section{پیشنهادها برای پژوهش‌های آتی}\label{sec:6-3}

بخش پایانی فصل به ارائه‌ی پیشنهادهای عملیاتی و راهبردی برای پژوهشگرانی که مایل به ادامه‌ی این مسیر هستند اختصاص می‌یابد. پیشنهادها باید مشخص، دست‌یافتنی و برآمده از یافته‌ها یا محدودیت‌های شناسایی‌شده در این پایان‌نامه باشند:
\begin{itemize}
    \item گسترش دامنه‌ی روش به سیستم‌های پیچیده‌تر با شرایط ناهمگن؛
    \vspace{-0.25cm}
    \item کاهش بیشتر پیچیدگی‌های محاسباتی با به‌کارگیری تقریب‌های عددی کارآمدتر؛
    \vspace{-0.25cm}
    \item پیاده‌سازی و ارزیابی تجربی روش در بسترهای عملیاتی و کاربردهای بلادرنگ.
\end{itemize}
